% RECOVERY32 reconstruction from the restored cycle31 proof.
% Not asserted byte-identical to the lost pre-outage cycle32 source.
\subsection*{The actual family-indexed basis of the relative fiber}
\statement{B.156}{Proved}{The paired relative fiber has a complete ordinary free-product basis retaining every deleted family core}
\label{stmt:B156}
In B.155 use the actual generic orbifold tuple
\[
 E=F(a,b,d,e,f)*C_2(C)*C_2(\tau),\qquad
 x=\tau C,\qquad z^*=\tau z\tau.
\]
The ordinary fiber has the literal decomposition
\begin{equation}\label{eq:B156-factors}
 F_{11}=R*D,\qquad
 R=F(x,a,a^*,f,f^*),\quad
 D=F(b,b^*,d,d^*,e,e^*).
\end{equation}
Let $N_a=\ker(R\to W_a)$ and $R_U=(R\to W_a)^{-1}(U)$.
For a right-coset transversal with identity representative $1$,
\begin{equation}\label{eq:B156-star}
 Q=R_U*
 \left(*_{R_Ut\in R_U\backslash R}tDt^{-1}\right).
\end{equation}
Representatives can be chosen in exact $\delta$ pairs except for $1$.
The displayed six deleted-family cores and all of their coset copies
are therefore simultaneous free-basis letters exchanged in pairs.
The retained factor $R_U$ also admits a paired free basis.
These are ordinary statements; no compact free-product decomposition
of this infinite-index subgroup is asserted.
\begin{proof}
The two-coset Schreier rewriting of $E$ with representatives $1,\tau$
gives \eqref{eq:B156-factors}. The involution inverts $x$ and exchanges
each starred pair, so both free factors are stable. The deletion
$R*D\to R$ is the actual one-family deletion, and the filled one-family
fiber is $W_a=R/N_a$ with its complete projected node and source
closure. The original source actions on $F_{11}$ need not preserve
the literal factor $R$; the equivariant map is the deletion map.

The map $F_{11}\to W_a$ factors through this retraction. Therefore $Q$
is exactly the inverse image of $R_U$ under $R*D\to R$.
In the ordinary Bass--Serre tree, the quotient by $Q$ has one
$R$-type vertex, with stabilizer $R_U$, and one $D$-type leaf for
each coset $R_Ut$. Its edges have trivial stabilizers and its
underlying graph is a star. Its fundamental-group presentation is
exactly \eqref{eq:B156-star}. In particular no additional graph
free factor or extra relation is omitted. The maps are the literal
inclusions of $R_U$ and $tDt^{-1}$; ordinary reduced-word normal
form proves their inverse presentation as well.

The cosets $R_U\backslash R$ identify with $U\backslash W_a$.
Only the identity coset is $\delta$-fixed. Indeed if
$\delta(w)=uw$ with $u\in U$, applying $\delta$ twice gives $u^2=1$;
$U$ is torsion-free and is the fixed subgroup of $W_a$, so the coset
is the identity. Choose representatives of the remaining two-element
orbits and their exact $\delta$ images. On the identity $D$ factor
the involution exchanges $b,b^*$, $d,d^*$ and $e,e^*$.
On any other factor it sends
$tzt^{-1}$ to $\delta(t)z^*\delta(t)^{-1}$.
This gives the asserted simultaneous paired deleted-core basis.

For $R_U$, apply the one-cone argument of B.155 to the smaller literal
orbifold $C_2(C)*C_2(\tau)*F(a,f)$. The inverse image of
$U\times\langle\tau\rangle$ contains exactly one $\tau$ factor and no
$C$ factor. Choosing the remaining free basis to be even and taking
the ordinary kernel gives a free basis of $R_U$ exchanged in pairs.
Combining it with \eqref{eq:B156-star} proves the complete paired basis.

The actual final relation subgroup is still the FULL $N$ of B.155,
not merely the normal closure of the projected $N_a$. Thus
\[
 Q/N=
 \left(R_U*F(\text{all displayed deleted-core copies})\right)/N,
 \qquad Q/N\longrightarrow U,\quad\ker=K_{\mathrm{rem}}.
\]
Every word from B.155's complete node scheme is rewritten with its
actual connectors and signed source corrections. Although
$R_U/N_a=U$, the map $R_U\to Q/N$ is not asserted to factor through
$U$: a projected relator may still represent a nontrivial remaining
kernel word before the other family letters are deleted. The actual
section $\sigma(U)$ of B.149/B.154 need not be this literal free factor.

Finally $N_a\le[R,R]$, because the entire projected nodal family
consists of commutators and all its translates remain commutators.
Hence $(W_a)_{\mathrm{ab}}=R_{\mathrm{ab}}$. The source actions
preserve the $a,a^*$ peripheral abelian classes, and $U$ is generated
by paired norms of their conjugates. Its abelian image lies on the
diagonal line generated by $[a]+[a^*]$. A retained core $a$ or $a^*$,
with any actual conjugator, has one of the two separate unit-vector
classes and therefore does not belong to $R_U$ or $Q$.
Thus a node commutator lying in $Q$ cannot be treated as a commutator
of two $Q$-basis vertices merely because its full word lies there.
All of its Schreier path and connector must be kept.
\end{proof}
The complete proof is
\path{work_CE/remaining_kernel_family_basis31.txt}.
All original nodes, both sheets, signed operators and retained-deck
translates remain. No filled-RAAG or remaining fixed-generation
conclusion follows from the basis decomposition alone.
